\newglossaryentry{dtheorem}{
  name = {演绎定理},
  description = {将基于假设的语句的语义后承和可证性，与对应的条件句的语义后承和可证性联系起来的定理。
在语义形式中（\olref[fol][syn][sem]{thm:sem-deduction}），它断言 $\Gamma \cup \{!A\} \Entails !B$ 当且仅当 $\Gamma \Entails !A \lif
    !B$。在证明论形式中，它断言 $\Gamma \cup
   \{!A\} \Proves !B$ 当且仅当 $\Gamma \Proves !A \lif !B$}}

\newglossaryentry{set-extensionality} {
  name={外延性（集合的）},
  description={集合 $A$ 和~$B$ 相等，即 $A = B$，当且仅当~$A$ 的每个
    元素都是~$B$ 的元素，反之亦然
    （参见 \olref[sfr][set][bas]{sec}）}}

\newglossaryentry{sat-extensionality}{
        name = {外延性（满足关系的）},
        description = {公式~$!A$ 是否被满足，
        仅取决于实际出现在~$!A$ 中的非逻辑符号和自由变元的赋值}}

\newglossaryentry{set} {
  name=集合,
  description={对象的汇集，独立于其指定方式、集合中对象的顺序
    以及对象的多重性（参见 \olref[sfr][set][bas]{sec}）} }

\newglossaryentry{symmetric} {
  name=对称的,
  description={关系 $R$ 是对称的，当且仅当只要 $Rxy$ 成立，就有 $Ryx$ 成立
    （参见 \olref[sfr][rel][prp]{sec}）} }

\newglossaryentry{cp} {
  name=笛卡尔积,
  symbol = {\ensuremath{A \times B}},
  description={所有由 $A$ 和 $B$ 的{元素}组成的有序对的集合；$A
    \times B = \Setabs{\tuple{x, y}}{x \in A \text{ 且 } y \in B}$
    （参见 \olref[sfr][set][pai]{sec}）} }

\newglossaryentry{subset} {
  name=子集,
  symbol = {\ensuremath{A \subseteq B}},
  description = {一个集合，其每个{元素}都是某个给定集合~$B$ 的{元素}
    （参见 \olref[sfr][set][sub]{sec}）}  }

\newglossaryentry{finsequence} {
  name = {序列（有限的）},
  symbol = {\ensuremath{A^*}},
  description = {由~$A$ 的{元素}组成的有穷字符串；是某个 $n$ 对应的 $A^n$ 的
    {元素}（参见
    \olref[sfr][set][imp]{sec}）}}

\newglossaryentry{infsequence} {
  name = {序列（无穷的） },
  symbol = {\ensuremath{A^\omega}},
  description = {一个没有间隙、永不终止的、由~$A$ 的{元素}组成的序列；
    形式上，是一个函数 $s\colon \PosInt \to A$（参见
    \olref[sfr][set][imp]{sec}）}}

\newglossaryentry{ps} {
  name=幂集,
  symbol = {\ensuremath{\Pow{A}}},
  description={由一个集合~$A$ 的所有\glspl{subset}组成的集合，
    $\Pow{A} = \Setabs{x}{x \subseteq A}$（参见
    \olref[sfr][set][sub]{sec}）} }

\newglossaryentry{union} {
  name=并集,
  symbol = {\ensuremath{A \cup B}},
  description={所有属于 $A$ 和 $B$ 的{元素}放在一起构成的集合：$A
    \cup B = \Setabs{x}{x \in A \lor x\in B}$（参见
    \olref[sfr][set][uni]{sec}）} }

\newglossaryentry{intersection} {
  name=交集,
  symbol = {\ensuremath{A \cap B}},
  description={所有同时是 $A$ 和~$B$ 的{元素}的事物的集合：$A \cap B = \Setabs{x}{x \in A \land x\in B}$（参见
    \olref[sfr][set][uni]{sec}）} }

\newglossaryentry{difference} {
  name=差集,
  symbol = {\ensuremath{A \setminus B}},
  description={所有属于 $A$ 但不属于~$B$ 的{元素}的集合：$A \setminus B = \Setabs{x}{x \in A \text{ 且 }
    x \notin B}$（参见 \olref[sfr][set][uni]{sec}）} }

\newglossaryentry{disjoint}{
  name=不相交,
  description={没有公共{元素}的两个集合（参见
    \olref[sfr][set][uni]{sec}）}}

\newglossaryentry{composition} {
  name=复合,
  symbol = {\ensuremath{g \circ f}},
  description={通过将 $f$ 和 $g$“链接”起来得到的函数；
    $(g \circ f)(x) = g(f(x))$（参见
    \olref[sfr][fun][cmp]{sec}）} }

\newglossaryentry{function} {
  name=函数,
  symbol = {\ensuremath{f \colon A \to B}},
  description={将\gls{fdomain}~$A$ 的每个元素映射到陪域~$B$ 的一个元素
    （参见 \olref[sfr][fun][bas]{sec}）} }

\newglossaryentry{fdomain} {
  name = {定义域（函数的）},
  symbol = {\ensuremath{\dom{f}}},
  description = {一个（偏）函数有定义的对象所组成的集合
    （参见 \olref[sfr][fun][bas]{sec}）}}

\newglossaryentry{range}{
  name = 值域,
  symbol = {\ensuremath{\ran{f}}},
  description = {陪域中实际被 $f$ 输出的那些元素构成的子集；$\ran{f} = \Setabs{y \in B}{f(x) = y \text{ 对某个 $x \in
        A$}}$（参见 \olref[sfr][fun][bas]{sec})}}

\newglossaryentry{graph} {
  name=图像（函数的）,
  description={由 $R_f
    = \Setabs{\tuple{x,y}}{f(x) = y}$ 定义的关系 $R_f \subseteq A \times B$，其中 $f\colon A \pto B$（参见
    \olref[sfr][fun][rel]{sec}）} }

\newglossaryentry{inverse} {
  name=逆函数,
  description={若 $f\colon A \to B$ 是一个\gls{bijection}，则 $f^{-1} \colon
    B \to A$ 是一个函数，满足 $f^{-1}(y) =$ 使得 $f(x) = y$ 成立的那个唯一 $x \in
    A$（参见 \olref[sfr][fun][inv]{sec}）} }

\newglossaryentry{pf} {
  name=偏函数,
    symbol = {\ensuremath{f \colon A \pto B}},
  description={偏函数是一个映射，它给~$A$ 的每个{元素}至多指派一个~$B$ 的{元素}。如果 $f$ 将 $B$ 中的一个元素指派给 $x \in A$，则 $f(x)$ 是有定义的，否则是无定义的（参见 \olref[sfr][fun][par]{sec}）} }

\newglossaryentry{tc} {
  name=传递闭包,
  symbol = {\ensuremath{R^+}},
  description={包含~$R$ 的最小\gls{transitive}关系（参见 \olref[sfr][rel][ops]{sec}）} }

\newglossaryentry{ir} {
  name=逆关系,
  symbol = {\ensuremath{R^{-1}}},  
  description={将关系 $R$ “反过来”；$R^{-1} =
    \Setabs{\tuple{y, x}}{\tuple{x, y} \in R}$（参见
    \olref[sfr][rel][ops]{sec}）} }

\newglossaryentry{preorder} {
  name=预序,
  description={一个\gls{reflexive}且\gls{transitive}的关系（参见
    \olref[sfr][rel][ord]{sec}）} }

\newglossaryentry{po} {
  name=偏序,
  description={一个\gls{reflexive}、\gls{anti-symmetric}、
    \gls{transitive}的关系（参见 \olref[sfr][rel][ord]{sec}）} }

\newglossaryentry{to} {
  name=全序,
  description={参见 \gls{lo}}}

\newglossaryentry{irreflexive} {
  name=反自反的,
  description={关系 $R$ 是反自反的，如果没有 $x \in A$ 使得 $Rxx$ 成立（参见
    \olref[sfr][rel][ord]{sec}）} }

\newglossaryentry{asymmetric} {
  name=非对称的,
  description={关系 $R$ 是非对称的，如果没有元素对 $x,y\in A$ 使得我们既有
    $Rxy$ 又有 $Ryx$（参见 \olref[sfr][rel][ord]{sec}）} }

\newglossaryentry{so} {
  name=严格序,
  description={一个\gls{irreflexive}、\gls{asymmetric}且
    \gls{transitive}的关系（参见 \olref[sfr][rel][ord]{sec}）} }

\newglossaryentry{slo} {
  name=严格线性序,
  description={一个\gls{connected}的\gls{so}（参见
    \olref[sfr][rel][ord]{sec}）} }

\newglossaryentry{lo} {
  name=线性序,
  description={一个\gls{connected}的\gls{po}（参见
    \olref[sfr][rel][ord]{sec}）} }

\newglossaryentry{br} {
  name=二元关系,
  description={集合~$A^{2}$ 的一个子集；我们用 $Rxy$（或 $xRy$）表示
    $\tuple{x, y} \in R$（参见 \olref[sfr][rel][set]{sec}）} }

\newglossaryentry{reflexive} {
  name=自反的,
  description={关系 $R$ 是自反的，当且仅当对每个 $x \in A$，有 $Rxx$（参见
    \olref[sfr][rel][prp]{sec}）} }

\newglossaryentry{transitive} {
  name=传递的,
  description={关系 $R$ 是传递的，当且仅当只要 $Rxy$ 且 $Ryz$，则
    必有 $Rxz$（参见 \olref[sfr][rel][prp]{sec}）} }

\newglossaryentry{anti-symmetric} {
  name=反对称的,
  description={关系 $R$ 是反对称的，当且仅当只要同时有 $Rxy$ 和
    $Ryx$，则 $x=y$；换言之：若 $x\neq y$，则或者没有 $Rxy$，或者没有 $Ryx$（参见 \olref[sfr][rel][prp]{sec}）} }

\newglossaryentry{connected} {
  name=连通的,
  description={关系 $R$ 是连通的，如果对所有 $x, y\in A$，
    只要 $x \neq y$，则要么 $Rxy$，要么~$Ryx$（参见
    \olref[sfr][rel][prp]{sec}）} }

\newglossaryentry{er} {
  name={等价关系}, 
  description={一个自反、对称且传递的关系（参见
    \olref[sfr][rel][prp]{sec}）} }

\newglossaryentry{sf} {
  name={满射的},
  description={$f \colon A \to B$ 是满射的当且仅当~$f$ 的
    值域就是整个~$B$，即，对每个 $y \in B$ 至少存在一个 $x \in A$ 使得~$f(x) = y$（参见
    \olref[sfr][fun][kin]{sec}）} }

\newglossaryentry{if} {
  name={单射的}, 
  description={$f \colon A \to B$ 是单射的当且仅当对
    每个 $y \in B$，至多存在一个 $x \in A$ 使得~$f(x) =
    y$；等价地说，若 $x \neq x'$ 则 $f(x) \neq f(x')$
    （参见 \olref[sfr][fun][kin]{sec}）} }

\newglossaryentry{bijection} {
  name={双射},  
  description={一个既是满射又是单射的函数（参见 \olref[sfr][fun][kin]{sec}）} }

\newglossaryentry{enumeration} {
  name=枚举,  
  description={一个集合~$A$ 的所有元素的可能无限的列表；形式化地说，是一个满射函数
    $f\colon \Nat \to A$（参见 \olref[sfr][siz][enm]{sec}）} }

\newglossaryentry{equinumerous} {
  name=等势的,
  description={$A$ 和 $B$ 是等势的当且仅当存在一个从 $A$ 到 $B$ 的全
    双射（参见 \olref[sfr][siz][equ]{sec}）} }

\newglossaryentry{finitely satisfiable} {
  name={有限可满足的},  
  description={$\Gamma$ 是有限可满足的当且仅当每个有限
    $\Gamma_0 \subseteq \Gamma$ 都是可满足的（参见
    \olref[fol][com][com]{sec}）} }

\newglossaryentry{mcs} {
  name=完备一致集,
  description={一个\gls{sentence}集合是完备且一致的，当且仅当它
    是一致的，并且对每个\gls{sentence}~$!A$，要么 $!A$ 要么
    $\lnot !A$ 在集合中（参见 \olref[fol][com][ccs]{sec}）} }

\newglossaryentry{mos} {
  name=模型,  
  description={一个在其中每个\gls{sentence}
    ~$\Gamma$ 中的语句都为真的\gls{structure}就是~$\Gamma$ 的模型（参见
    \olref[fol][mat][exs]{sec}）} }

\newglossaryentry{closed} {
  name=封闭的, 
  description={一个\gls{sentence}集合~$\Gamma$ 是封闭的，当且仅当只要
    $\Gamma \Entails !A$ 则 $!A \in \Gamma$。集合
    $\Setabs{!A}{\Gamma \Entails !A}$ 是~$\Gamma$ 的闭包（参见
    \olref[fol][mat][int]{sec}）} }

\newglossaryentry{derivability} {
  name=可推导性,
  symbol = {\ensuremath{\Gamma \Proves !A}},
  description={在相继式演算中，如果存在一个\gls{sequent} $\Gamma_0 \Sequent !A$ 的\gls{derivation}，其中 $\Gamma_0 \subseteq \Gamma$ 是~$\Gamma$ 中语句的有穷序列，则称 $!A$ 从~$\Gamma$ 是{可推导的}（参见
    \olref[fol][seq][ptn]{sec}）。
在自然演绎中，如果存在一个以 $!A$ 为结束\gls{formula}、并且其中每个假设要么被{解除}要么属于~$\Gamma$ 的{推导}，则称 $!A$ 从~$\Gamma$ 是{可推导的}（参见
    \olref[fol][ntd][ptn]{sec}）} }

\newglossaryentry{theorem} {
  name={定理},
  symbol = {\ensuremath{\Proves !A}},
  description={在相继式演算中，如果存在相继式 $\Sequent !A$ 的{推导}，则称{公式}~$!A$ 是（逻辑的）{定理}（参见 \olref[fol][seq][ptn]{sec}）。
在自然演绎中，如果存在一个所有假设均被\gls{discharged}的 $!A$ 的{推导}，则称{公式}~$!A$ 是{定理}（参见 \olref[fol][ntd][ptn]{sec}）。
我们也称 $!A$ 是理论~$\Gamma$ 的定理，如果 $\Gamma \Proves !A$} }

\newglossaryentry{consistent} {
  name=一致的,
  description={在相继式演算中，语句集~$\Gamma$ 是{一致的}，当且仅当不存在以 $\Gamma_0 \subseteq \Gamma$ 为前提的\gls{sequent} $\Gamma_0 \Sequent \quad$ 的\gls{derivation}（参见 \olref[fol][seq][ptn]{sec}）。
在自然演绎中，$\Gamma$ 是{一致的}，当且仅当 $\Gamma \Proves/ \lfalse$（参见 \olref[fol][ntd][ptn]{sec}）。
如果 $\Gamma$ 不是一致的，则它是{不一致的}。} }

\newglossaryentry{inconsistent} {
  name=不一致的,
  description = {参见\gls{consistent}}}

\newglossaryentry{derivation} {
  name=推导,
  description={在相继式演算中，指一个相继式的树，其中每个相继式或者是初始相继式，或者由紧接其上方的相继式通过某条推理规则得到（参见 \olref[fol][seq][rul]{sec}）。
在自然演绎中，指一个公式的树，其中每个公式或者是假设，或者由紧接其上方的公式通过某条推理规则得到（参见 \olref[fol][ntd][rul]{sec}）} }

\newglossaryentry{eigenvariable}{
  name = 特征变元,
  description = {在相继式演算中，指 $\LeftR{\lexists}$ 或 $\RightR{\lforall}$ 推理的前提中出现、但不得在结论中出现的一个特殊常项符号（参见
    \olref[fol][seq][rul]{sec}）。
在自然演绎中，指 $\Elim{\lexists}$ 或 $\Intro{\lforall}$ 推理的前提中出现、但不得在结论或任何\gls{undischarged}假设中出现的一个特殊常项符号（参见
    \olref[fol][ntd][rul]{sec}）}}

\newglossaryentry{discharged}{
  name = {已解除的},
  description = {推导中的一个\gls{assumption}可以被其下方的某条推理规则{解除}（该规则与该假设会分配一个匹配的标签，例如 $[!A]^2$）。
如果它未被解除，则称为{未解除的}（参见 \olref[fol][ntd][rul]{sec}）}}

\newglossaryentry{undischarged}{
  name = {未解除的},
  description = {参见\gls{discharged}}}

\newglossaryentry{assumption}{
  name = {假设},
  description = {位于推导最顶端的公式，也称为初始公式。它可以是\gls{discharged}或未解除的（参见\olref[fol][ntd][rul]{sec}）}}

\newglossaryentry{soundness} {
  name=可靠性,
  description={推导系统的性质：若每当
    $\Gamma \Proves !A$ 时都有 $\Gamma \Entails !A$，则它是可靠的（参见
    \olref[fol][seq][sou]{sec} 和 \olref[fol][ntd][sou]{sec}）} }

\newglossaryentry{sequent}{
  name=矢列,
  description = {形如 $\Gamma \Sequent \Delta$ 的表达式，其中 $\Gamma$ 和 $\Delta$ 是语句的有限序列（参见 \olref[fol][seq][rul]{sec}）} }

\newglossaryentry{sentence} {
  name=语句,
  description={不含\gls{free}变元的\gls{formula}（参见
    \olref[fol][syn][fvs]{sec}）} }

\newglossaryentry{variable assignment} {
  name=变元指派,
  description={将每个{变元}映射到~$\Domain M$ 中一个元素的 \gls{function}（参见 \olref[fol][syn][sat]{sec}）} }

\newglossaryentry{$x$-variant} {
  name=$x$-变体,
  sort={x-variant},
  description={若两个\glspl{variable assignment}至多在赋予~$x$ 的值上不同，则称它们是 $x$-变体，记作 $s \sim_x
    s'$（参见
    \olref[fol][syn][sat]{sec}）} }

\newglossaryentry{valid} {
  name=有效,
  symbol = {\ensuremath{\Entails !A}},
  description={一个语句 $!A$ 是\emph{有效的}，当且仅当
    对每个 \gls{structure}~$\Struct M$ 都有 $\Sat{M}{!A}$（参见
    \olref[fol][syn][sem]{sec}）} }

\newglossaryentry{entailment} {
  name=语义后承,
  symbol = {\ensuremath{\Gamma \Entails !A}},
  description={一组语句~$\Gamma$ \emph{语义后承}一个
    语句~$!A$，当且仅当对每个满足 $\Sat{M}{\Gamma}$ 的
    \gls{structure}~$\Struct M$，都有 $\Sat{M}{!A}$
    （参见 \olref[fol][syn][sem]{sec}）} }

\newglossaryentry{satisfiable} {
  name=可满足,
  description={一组语句~$\Gamma$ 是\emph{可满足的}，若
    存在 \gls{structure}~$\Struct M$ 使得 $\Sat{M}{\Gamma}$，否则它是
    不可满足的（参见 \olref[fol][syn][sem]{sec}）} }

\newglossaryentry{cs} {
    name=覆盖结构,
    description={一个 \gls{structure}，其论域中每个元素都是某个闭项的值（参见
      \olref[fol][syn][str]{sec}）} }

\newglossaryentry{subformula} {
  name={子公式},
  description={公式的一部分，且本身是一个公式（参见
    \olref[fol][syn][sbf]{sec}）} }

\newglossaryentry{free for} {
  name=对……自由,
  description={项~$t$ \emph{对} $x$ \emph{自由}于 $!A$ 中，如果在 $!A$ 中~$x$ 的所有
    \gls{free} 出现都不在能够绑定~$t$ 中变元的量词的辖域内（参见
    \olref[fol][syn][sub]{sec}）} }

\newglossaryentry{church-turing thesis} {
  name=Church-Turing论题,
  description={断言：凡是通过有效程序可计算的，都是 Turing可计算的（参见 \olref[tur][mac][ctt]{sec}）} }

\newglossaryentry{completeness} {
  name=完备性,
  description={推导系统的性质；若每当
    $\Gamma$ 语义上后承 $!A$ 时，都存在一个\gls{derivation}可建立 $\Gamma \Proves !A$，则该系统是完备的；等价地说，当且仅当
    每个 \gls{consistent} 语句集都是 \gls{satisfiable} 的
    （参见 \olref[fol][com][int]{sec}）} }

\newglossaryentry{formula} {
  name={公式},
  description={一阶语言~$\Lang L$ 中表达关系或性质，或具有真假的表达式（参见
    \olref[fol][syn][frm]{sec}）} }

\newglossaryentry{structure} {
  name={结构},
  symbol = {\ensuremath{\Struct{M}}},
  description={对一阶语言的解释，由\gls{sdomain}以及对语言的常项、
    谓词和函数符号的指派组成（参见
    \olref[fol][syn][str]{sec}）} }

\newglossaryentry{sdomain} {
  name = {论域（结构的）},
  symbol = {\ensuremath{\Domain{M}}},
  description = {非空集合，\gls{structure} 从中获取指派和变元的值（参见
    \olref[fol][syn][str]{sec}）} }

\newglossaryentry{halting problem} {
  name=停机问题,
  description={判定（对任意 $e$，$n$） Turing机~$M_e$ 在输入~$n$ 个划痕时是否停机的问题（参见
    \olref[tur][und][hal]{sec}）} }

\newglossaryentry{bound} {
  name=约束出现,
  description={在同一变元的量词辖域内出现的该变元的出现
    （参见 \olref[fol][syn][fvs]{sec}）}
}

\newglossaryentry{free} {
  name=自由出现,
    description={不是\gls{bound}出现的变元出现
    （参见 \olref[fol][syn][fvs]{sec}）}
}

\newglossaryentry{compactness theorem} {
  name = {紧致性定理},
  description = {断言：每个 \gls{finitely satisfiable} 语句集都是可满足的（参见
    \olref[fol][com][com]{sec}）}}

\newglossaryentry{ls}{
  name = {Löwenheim-Skolem定理},
  description = {断言：每个可满足的语句集都有一个可数模型（参见 \olref[fol][com][dls]{sec}）}}

\newglossaryentry{Church-Turing Theorem}{
  name = {Church-Turing定理},
  description = {断言：不存在能判定任意给定一阶逻辑语句是否\gls{valid}的 Turing机（参见
    \olref[tur][und][uns]{sec}）}}

\newglossaryentry{decision problem}{
  name = {判定问题},
  description = {判定给定一阶逻辑语句是否\gls{valid}的问题（参见 \gls{Church-Turing
    Theorem}）}}

\newglossaryentry{completeness-theorem}{
  name = {完备性定理},
  description = {断言一阶逻辑是完备的：每个
    \gls{consistent} 语句集都是 \gls{satisfiable} 的}}
